%========================================
\chapter{应用层}
%========================================
\section{应用层概述}
应用层是TCP/IP协议栈的最高层，直接为用户提供网络服务。

\subsection{客户-服务器模式}
\begin{itemize}[leftmargin=2em,itemsep=0pt]
    \item \textbf{服务器}：提供服务，被动等待请求，通常使用固定端口
    \item \textbf{客户端}：请求服务，主动发起连接，使用临时端口
\end{itemize}

\subsection{P2P模式}
\begin{itemize}[leftmargin=2em,itemsep=0pt]
    \item 每个节点既是客户端又是服务器
    \item 资源分布在各个节点上
    \item 可扩展性好，但管理复杂
\end{itemize}

\section{DNS域名系统}
DNS（Domain Name System）将域名解析为IP地址。

\subsection{域名结构}
域名采用层次结构，从右到左依次为：根域、顶级域、二级域、三级域...

\textbf{示例}：\texttt{www.example.com}
\begin{itemize}[leftmargin=2em,itemsep=0pt]
    \item \texttt{com}：顶级域（通用顶级域）
    \item \texttt{example}：二级域
    \item \texttt{www}：主机名
\end{itemize}

\subsection{DNS服务器类型}
\begin{itemize}[leftmargin=2em,itemsep=0pt]
    \item \textbf{根域名服务器}：知道所有顶级域名服务器的地址
    \item \textbf{顶级域名服务器}：管理该顶级域下的所有二级域名
    \item \textbf{权威域名服务器}：负责特定区域的域名解析
    \item \textbf{本地域名服务器}：用户首先查询的服务器，有缓存功能
\end{itemize}

\subsection{DNS查询方式}
\begin{itemize}[leftmargin=2em,itemsep=0pt]
    \item \textbf{递归查询}：本地DNS服务器代替客户端完成全部查询
    \item \textbf{迭代查询}：本地DNS服务器依次向各级服务器查询
\end{itemize}

\begin{center}
\begin{tikzpicture}[>=stealth,thick,scale=0.85]
\node[draw,rectangle] (client) at (0,0) {客户端};
\node[draw,rectangle] (local) at (3,0) {本地DNS};
\node[draw,rectangle] (root) at (6,2) {根DNS};
\node[draw,rectangle] (tld) at (6,0) {顶级DNS};
\node[draw,rectangle] (auth) at (6,-2) {权威DNS};
\draw[->] (client) -- node[above]{\tiny 1.查询} (local);
\draw[->] (local) -- node[above,sloped]{\tiny 2} (root);
\draw[->] (root) -- node[right]{\tiny 3} (local);
\draw[->] (local) -- node[above]{\tiny 4} (tld);
\draw[->] (tld) -- node[right]{\tiny 5} (local);
\draw[->] (local) -- node[above,sloped]{\tiny 6} (auth);
\draw[->] (auth) -- node[right]{\tiny 7} (local);
\draw[->] (local) -- node[below]{\tiny 8.结果} (client);
\end{tikzpicture}
\end{center}

\begin{notebox}
\textbf{DNS使用的协议}：
\begin{itemize}[itemsep=0pt]
    \item 通常使用UDP（端口53），查询响应快
    \item 当响应超过512字节时使用TCP
    \item 区域传送使用TCP
\end{itemize}
\end{notebox}

\begin{examplebox}[例题6.1 DNS查询]
\textbf{题目}：用户访问www.example.com，本地DNS缓存中没有记录，描述DNS解析过程。
\tcblower
\begin{enumerate}[itemsep=0pt]
    \item 客户端向本地DNS服务器发送查询请求
    \item 本地DNS向根服务器查询，获得.com顶级域服务器地址
    \item 本地DNS向.com服务器查询，获得example.com权威服务器地址
    \item 本地DNS向权威服务器查询，获得www.example.com的IP地址
    \item 本地DNS将结果返回给客户端，并缓存该记录
\end{enumerate}
\end{examplebox}

\section{HTTP协议}
HTTP（HyperText Transfer Protocol）是Web应用的基础协议。

\subsection{HTTP特点}
\begin{itemize}[leftmargin=2em,itemsep=0pt]
    \item \textbf{无状态}：服务器不保存客户端状态信息
    \item \textbf{基于TCP}：使用TCP连接，默认端口80
    \item \textbf{请求-响应模式}：客户端发送请求，服务器返回响应
\end{itemize}

\subsection{HTTP连接方式}
\begin{itemize}[leftmargin=2em,itemsep=0pt]
    \item \textbf{非持久连接}（HTTP/1.0）：每个请求/响应使用单独的TCP连接
    \item \textbf{持久连接}（HTTP/1.1）：多个请求/响应复用同一TCP连接
    \begin{itemize}[itemsep=0pt]
        \item 非流水线：收到响应后才发送下一个请求
        \item 流水线：连续发送多个请求，不等待响应
    \end{itemize}
\end{itemize}

\subsection{HTTP请求方法}
\begin{itemize}[leftmargin=2em,itemsep=0pt]
    \item \textbf{GET}：获取资源
    \item \textbf{POST}：提交数据
    \item \textbf{PUT}：上传文件
    \item \textbf{DELETE}：删除资源
    \item \textbf{HEAD}：获取响应头（不返回内容）
\end{itemize}

\subsection{HTTP状态码}
\begin{table}[htbp]
\centering\small
\begin{tabular}{cl}
\toprule
\textbf{状态码} & \textbf{含义} \\
\midrule
1xx & 信息性状态码 \\
2xx & 成功（200 OK） \\
3xx & 重定向（301永久移动，302临时移动） \\
4xx & 客户端错误（404未找到，403禁止） \\
5xx & 服务器错误（500内部错误） \\
\bottomrule
\end{tabular}
\end{table}

\begin{examplebox}[例题6.2 HTTP时延计算]
\textbf{题目}：用户访问一个包含1个HTML文件和2个图片的网页，RTT=100ms。分别计算非持久连接和持久连接（非流水线）的总时延。
\tcblower
\textbf{非持久连接}：每个对象需要2RTT（建立连接+请求响应）

总时延 = 3 × 2RTT = 6RTT = \textbf{600ms}

\textbf{持久连接（非流水线）}：第一个对象2RTT，后续每个1RTT

总时延 = 2RTT + 2 × 1RTT = 4RTT = \textbf{400ms}

\textbf{持久连接（流水线）}：建立连接后可连续请求

总时延 ≈ 2RTT + 1RTT = 3RTT = \textbf{300ms}
\end{examplebox}

\subsection{Cookie与Session}
\begin{itemize}[leftmargin=2em,itemsep=0pt]
    \item \textbf{Session}：存储在服务器端，通过Session ID关联
\end{itemize}

\subsection{HTTPS 协议}
HTTPS (HyperText Transfer Protocol Secure) 是以安全为目标的 HTTP 通道，在 HTTP 下加入了 SSL/TLS 层。

\textbf{HTTPS 与 HTTP 的本质区别}：
\begin{itemize}[itemsep=0pt]
    \item \textbf{安全性}：HTTP 是明文传输，不安全；HTTPS 通过 SSL/TLS 进行加密传输、身份认证。
    \item \textbf{端口}：HTTP 默认使用 80 端口；HTTPS 默认使用 443 端口。
    \item \textbf{资源消耗}：HTTPS 由于需要加密解密，对 CPU 消耗更高，链接建立更慢。
    \item \textbf{证书}：HTTPS 需要 CA 机构颁发的 SSL 证书，通常需要费用。
\end{itemize}

\textbf{HTTPS 工作过程}：
\begin{enumerate}[leftmargin=2em,itemsep=0pt]
    \item \textbf{建立 TCP 连接}：进行标准的三次握手。
    \item \textbf{TLS 握手}：
    \begin{itemize}[itemsep=0pt]
        \item 客户端向服务器索要公钥（服务器证书）。
        \item 客户端验证证书的有效性及颁发机构。
        \item 双方协商生成“对话密钥”（对称加密密钥）。
    \end{itemize}
    \item \textbf{加密通信}：使用协商好的对称密钥加密应用层数据（HTTP）。
\end{enumerate}

\section{FTP文件传输协议}
FTP（File Transfer Protocol）用于在网络上传输文件。

\subsection{FTP特点}
\begin{itemize}[leftmargin=2em,itemsep=0pt]
    \item 使用两个TCP连接：
    \begin{itemize}[itemsep=0pt]
        \item \textbf{控制连接}（端口21）：传输命令和响应，持久连接
        \item \textbf{数据连接}（端口20）：传输文件数据，非持久连接
    \end{itemize}
    \item 支持ASCII和二进制传输模式
    \item 需要用户认证
\end{itemize}

\subsection{FTP工作模式}
\begin{itemize}[leftmargin=2em,itemsep=0pt]
    \item \textbf{主动模式}：服务器主动连接客户端的数据端口
    \item \textbf{被动模式}：客户端主动连接服务器的数据端口（适用于防火墙环境）
\end{itemize}

\section{电子邮件}
\subsection{邮件系统组成}
\begin{itemize}[leftmargin=2em,itemsep=0pt]
    \item \textbf{用户代理}：邮件客户端软件
    \item \textbf{邮件服务器}：存储和转发邮件
    \item \textbf{邮件协议}：SMTP、POP3、IMAP
\end{itemize}

\subsection{邮件协议}
\begin{itemize}[leftmargin=2em,itemsep=0pt]
    \item \textbf{SMTP}（端口25）：发送邮件，推送协议
    \item \textbf{POP3}（端口110）：接收邮件，下载后可删除服务器副本
    \item \textbf{IMAP}（端口143）：接收邮件，邮件保留在服务器
\end{itemize}

\begin{center}
\begin{tikzpicture}[>=stealth,thick,scale=0.9]
\node[draw,rectangle] (ua1) at (0,0) {发送方UA};
\node[draw,rectangle] (ms1) at (3,0) {发送方服务器};
\node[draw,rectangle] (ms2) at (6,0) {接收方服务器};
\node[draw,rectangle] (ua2) at (9,0) {接收方UA};
\draw[->] (ua1) -- node[above]{\tiny SMTP} (ms1);
\draw[->] (ms1) -- node[above]{\tiny SMTP} (ms2);
\draw[->] (ms2) -- node[above]{\tiny POP3/IMAP} (ua2);
\end{tikzpicture}
\end{center}

\section{DHCP动态主机配置协议}
DHCP（Dynamic Host Configuration Protocol）自动为主机分配IP地址等网络配置。

\subsection{DHCP工作过程}
\begin{enumerate}[leftmargin=2em,itemsep=0pt]
    \item \textbf{DHCP Discover}：客户端广播发现DHCP服务器
    \item \textbf{DHCP Offer}：服务器提供IP地址等配置
    \item \textbf{DHCP Request}：客户端请求使用某个配置
    \item \textbf{DHCP ACK}：服务器确认分配
\end{enumerate}

\section{SNMP 简单网络管理协议}
SNMP (Simple Network Management Protocol) 用于管理和监控网络设备。
\begin{itemize}[leftmargin=2em,itemsep=0pt]
    \item \textbf{架构}：基于 管理站 (Manager) 和 代理 (Agent) 模型。
    \item \textbf{工作方式}：使用 UDP（端口 161 为代理，162 为管理站 Traps）。
    \item \textbf{功能}：查询设备状态、修改配置、接收告警信息。
\end{itemize}

\begin{notebox}
\textbf{DHCP使用UDP}：
\begin{itemize}[itemsep=0pt]
    \item 客户端使用端口68
    \item 服务器使用端口67
    \item 初始时客户端没有IP地址，使用0.0.0.0作为源地址
\end{itemize}
\end{notebox}

\section{常用端口号总结}
\begin{table}[htbp]
\centering\small
\begin{tabular}{ccc}
\toprule
\textbf{协议} & \textbf{端口} & \textbf{传输层协议} \\
\midrule
FTP数据 & 20 & TCP \\
FTP控制 & 21 & TCP \\
SSH & 22 & TCP \\
Telnet & 23 & TCP \\
SMTP & 25 & TCP \\
DNS & 53 & UDP/TCP \\
DHCP服务器 & 67 & UDP \\
DHCP客户端 & 68 & UDP \\
HTTP & 80 & TCP \\
POP3 & 110 & TCP \\
IMAP & 143 & TCP \\
HTTPS & 443 & TCP \\
\bottomrule
\end{tabular}
\end{table}

\begin{examplebox}[例题6.3 应用层协议]
\textbf{题目}：下列协议中，使用UDP的是？\\
A. HTTP \quad B. FTP \quad C. SMTP \quad D. DNS
\tcblower
\begin{itemize}[itemsep=0pt]
    \item HTTP：TCP（端口80）
    \item FTP：TCP（端口20/21）
    \item SMTP：TCP（端口25）
    \item DNS：主要使用UDP（端口53），大数据量时用TCP
\end{itemize}
\textbf{答案}：D
\end{examplebox}

\begin{examplebox}[例题6.4 邮件协议]
\textbf{题目}：用户A使用Outlook发送邮件给用户B，B使用网页版Gmail查看邮件。描述邮件传输过程中使用的协议。
\tcblower
\begin{enumerate}[itemsep=0pt]
    \item A的Outlook通过\textbf{SMTP}将邮件发送到A的邮件服务器
    \item A的邮件服务器通过\textbf{SMTP}将邮件转发到Gmail服务器
    \item B通过浏览器访问Gmail，使用\textbf{HTTP/HTTPS}获取邮件
\end{enumerate}

注：如果B使用邮件客户端，则使用POP3或IMAP从服务器获取邮件。
\end{examplebox}

\begin{examplebox}[综合题6.3b HTTP请求分析]
\textbf{题目}：分析以下捕获到的 HTTP 请求报文。
\begin{lstlisting}[basicstyle=\ttfamily\footnotesize]
GET /index.html HTTP/1.1
Host: www.example.com
Connection: keep-alive
User-Agent: Mozilla/5.0 (Windows NT 10.0; Win64; x64)
Accept-Language: zh-CN,zh;q=0.9
\end{lstlisting}
\tcblower
\textbf{解析}：
\begin{enumerate}[itemsep=0pt]
    \item \textbf{请求行}：
    \begin{itemize}[itemsep=0pt]
        \item \textbf{方法}：GET，表示请求获取资源。
        \item \textbf{URL}：/index.html，请求的具体路径。
        \item \textbf{版本}：HTTP/1.1，目前主流的协议版本。
    \end{itemize}
    \item \textbf{Host}：指明目的服务器的域名为 \texttt{www.example.com}。
    \item \textbf{Connection}：keep-alive，表示使用持久连接。
    \item \textbf{User-Agent}：标识客户端浏览器的类型和操作系统版本。
    \item \textbf{Accept-Language}：说明客户端首选的语言是中文 (zh-CN)。
\end{enumerate}
\end{examplebox}

\section{Web缓存与代理服务器}
\subsection{Web缓存的作用}
\begin{itemize}[leftmargin=2em,itemsep=0pt]
    \item 减少客户端请求的响应时间
    \item 减少机构接入链路的流量
    \item 降低源服务器的负载
\end{itemize}

\subsection{条件GET}
用于验证缓存对象是否过期：
\begin{itemize}[leftmargin=2em,itemsep=0pt]
    \item 请求头包含\texttt{If-Modified-Since}字段
    \item 若对象未修改，服务器返回304 Not Modified（不含对象）
    \item 若对象已修改，服务器返回200 OK（含新对象）
\end{itemize}

\section{应用层综合例题}

\begin{examplebox}[例题6.5 DNS查询类型]
\textbf{题目}：本地DNS服务器向根DNS服务器的查询通常采用什么方式？为什么？
\tcblower
通常采用\textbf{迭代查询}。

原因：
\begin{itemize}[itemsep=0pt]
    \item 根DNS服务器负载很重，需要处理大量查询
    \item 迭代查询让本地DNS服务器承担主要工作
    \item 根服务器只需返回下一级服务器地址，减轻负担
\end{itemize}

而本地DNS服务器与客户端之间通常采用\textbf{递归查询}，因为客户端只需发送一次请求即可获得最终结果。
\end{examplebox}

\begin{examplebox}[例题6.6 HTTP持久连接]
\textbf{题目}：某网页包含1个HTML文件（1KB）和10个图片（每个10KB），RTT=100ms，链路带宽1Mbps。分别计算非持久连接和持久连接（流水线）方式下载该网页的总时间。
\tcblower
\textbf{传输时间计算}：
\begin{itemize}[itemsep=0pt]
    \item HTML传输时间 = $\frac{1\times8}{1000}$ = 8ms
    \item 每个图片传输时间 = $\frac{10\times8}{1000}$ = 80ms
\end{itemize}

\textbf{非持久连接}：每个对象需要2RTT + 传输时间
\begin{itemize}[itemsep=0pt]
    \item HTML：2×100 + 8 = 208ms
    \item 10个图片：10×(2×100 + 80) = 2800ms
    \item 总时间 = 208 + 2800 = \textbf{3008ms}
\end{itemize}

\textbf{持久连接（流水线）}：1次建立连接 + 连续请求
\begin{itemize}[itemsep=0pt]
    \item 建立连接：1RTT = 100ms
    \item 请求HTML并等待：1RTT + 8ms = 108ms
    \item 连续请求10个图片：1RTT + 10×80ms = 900ms
    \item 总时间 ≈ 100 + 108 + 900 = \textbf{1108ms}
\end{itemize}
\end{examplebox}

\begin{examplebox}[例题6.7 DHCP工作过程]
\textbf{题目}：主机A首次接入网络，描述其通过DHCP获取IP地址的过程。
\tcblower
\begin{enumerate}[itemsep=0pt]
    \item \textbf{DHCP Discover}：A广播发现报文
    \begin{itemize}[itemsep=0pt]
        \item 源IP：0.0.0.0（A还没有IP）
        \item 目的IP：255.255.255.255（广播）
        \item 源MAC：A的MAC地址
        \item 目的MAC：FF-FF-FF-FF-FF-FF
    \end{itemize}
    
    \item \textbf{DHCP Offer}：服务器提供配置
    \begin{itemize}[itemsep=0pt]
        \item 包含：可用IP、子网掩码、默认网关、DNS服务器、租约时间
    \end{itemize}
    
    \item \textbf{DHCP Request}：A请求使用该配置（广播）
    
    \item \textbf{DHCP ACK}：服务器确认分配
\end{enumerate}

此后A可以使用分配的IP地址，租约到期前需要续租。
\end{examplebox}

\begin{examplebox}[例题6.8 应用层协议综合]
\textbf{题目}：用户在浏览器输入\texttt{www.example.com}并按回车，描述涉及的应用层协议。
\tcblower
\begin{enumerate}[itemsep=0pt]
    \item \textbf{DNS}：将域名\texttt{www.example.com}解析为IP地址
    
    \item \textbf{HTTP/HTTPS}：
    \begin{itemize}[itemsep=0pt]
        \item 浏览器与Web服务器建立TCP连接
        \item 发送HTTP GET请求获取网页
        \item 服务器返回HTML文档
    \end{itemize}
    
    \item 若网页包含其他资源（图片、CSS、JS）：
    \begin{itemize}[itemsep=0pt]
        \item 可能再次进行DNS查询（不同域名）
        \item 发送额外的HTTP请求获取资源
    \end{itemize}
\end{enumerate}

底层还涉及：TCP（传输层）、IP（网络层）、ARP（获取网关MAC地址）
\end{examplebox}

\section{本章习题}

\begin{examplebox}[填空题]
\textbf{1.} 综合布线系统包括\underline{工作区子系统}、\underline{水平子系统}、\underline{管理子系统}、\underline{垂直干线子系统}、\underline{设备间子系统}、\underline{建筑群子系统}六个子系统。
\end{examplebox}

\begin{examplebox}[选择题]
\textbf{1.} 标准的URL组成：服务器类型、主机名和路径及（\textbf{C}）\\
A. 客户名 \quad B. 浏览器名 \quad C. 文件名 \quad D. 进程名

\textbf{2.} 协议与服务说法正确的是（\textbf{C}）\\
A. 协议是垂直的\\
B. 服务是水平的\\
C. 在协议的控制下，上层对下层进行调用，下层向上层提供服务\\
D. 同层两个实体间必须保持连接

\textit{解析}：协议是水平的（对等层之间），服务是垂直的（相邻层之间）。
\end{examplebox}

\begin{examplebox}[简答题：OSI/RM与TCP/IP对比]
\textbf{题目}：简要说明OSI/RM与TCP/IP的优缺点。
\tcblower
\textbf{OSI/RM}：
\begin{itemize}[itemsep=0pt]
    \item 优点：层次清晰，理论完善，各层功能明确
    \item 缺点：过于复杂，实现困难，市场推广失败
\end{itemize}

\textbf{TCP/IP}：
\begin{itemize}[itemsep=0pt]
    \item 优点：简单实用，得到广泛应用，是事实标准
    \item 缺点：层次划分不够清晰，网络接口层几乎为空
\end{itemize}
\end{examplebox}

\begin{examplebox}[简答题：差错发生层次]
\textbf{题目}：在OSI模型中，各层都有差错控制过程，指出以下每种差错发生在OSI的哪些层中？
\begin{enumerate}[itemsep=0pt]
    \item 噪声使传输信道上的一个0变成1或一个1变成0
    \item 一个分组被传送到错误的目的站
    \item 收到一个序号错误的帧
    \item 一台打印机正在打印，突然收到一个错误指令要打印头回到本行的开始位置
    \item 一个半双工的会话中，正在发送数据的用户突然接收到对方用户发来的数据
\end{enumerate}
\tcblower
\begin{enumerate}[itemsep=0pt]
    \item \textbf{物理层}：比特传输错误
    \item \textbf{网络层}：路由选择错误
    \item \textbf{数据链路层}：帧序号错误
    \item \textbf{应用层}：应用程序指令错误
    \item \textbf{会话层}：会话管理错误（半双工控制）
\end{enumerate}
\end{examplebox}

\begin{examplebox}[综合题：交换机自学习]
\textbf{题目}：如图所示3个交换机互联，假设：A向B发送了一帧，C向B发送了一帧，B向A发送了一帧，L向A发送了一帧，L向D发送了一帧。请分析：此时S1、S2和S3的交换表内容分别是什么？
\tcblower
\textbf{S1交换表}：
\begin{center}\small
\begin{tabular}{|c|c|}
\hline
MAC地址 & 端口 \\
\hline
A的MAC地址 & 连接A的端口 \\
C的MAC地址 & 连接C的端口 \\
L的MAC地址 & 连接S2的端口 \\
\hline
\end{tabular}
\end{center}

\textbf{S2交换表}：
\begin{center}\small
\begin{tabular}{|c|c|}
\hline
MAC地址 & 端口 \\
\hline
A的MAC地址 & 连接S1的端口 \\
B的MAC地址 & 连接B的端口 \\
C的MAC地址 & 连接S1的端口 \\
D的MAC地址 & 连接S3的端口 \\
L的MAC地址 & 连接S3的端口 \\
\hline
\end{tabular}
\end{center}

\textbf{S3交换表}：
\begin{center}\small
\begin{tabular}{|c|c|}
\hline
MAC地址 & 端口 \\
\hline
A的MAC地址 & 连接S2的端口 \\
D的MAC地址 & 连接D的端口 \\
L的MAC地址 & 连接L的端口 \\
\hline
\end{tabular}
\end{center}

\textbf{分析过程}：交换机通过自学习记录源MAC地址与入端口的映射，通过泛洪或明确转发处理目的地址。
\end{examplebox}

\begin{examplebox}[选择题（补充）]
\textbf{3.} 管理计算机通信的规则称为（\textbf{A}）\\
A. 协议 \quad B. 介质 \quad C. 服务 \quad D. 网络操作系统

\textbf{4.} 将一条物理信道按时间分成若干时间片轮换地给多个信号使用，这就是（\textbf{B}）\\
A. 频分多路复用 \quad B. 时分多路复用\\
C. 空分多路复用 \quad D. 频分与时分混合多路复用

\textbf{5.} 世界上很多国家都相继组建了自己国家的公用数据网，现有的公用数据网大都采用（\textbf{C}）\\
A. 电路交换方式 \quad B. 报文交换方式\\
C. 分组交换方式 \quad D. 空分交换方式
\end{examplebox}

\begin{examplebox}[计算题6.9 协议头开销]
\textbf{题目}：一个系统的协议结构有N层，应用程序产生M字节长的报文。网络软件在每层都加上H字节长的协议头。那么，网络带宽中有多大比率用于协议头信息的传输？
\tcblower
每个报文上附加的头字节的总数 = H × N

总传输数据量 = M + H × N

\textbf{协议头开销比率} = $\frac{H \times N}{M + H \times N}$

例如：若N=4，H=20，M=1000，则开销比率 = $\frac{80}{1080}$ ≈ 7.4\%
\end{examplebox}

\begin{examplebox}[计算题6.10 HTTP响应时间]
\textbf{题目}：某网页包含1个HTML文件（2KB）和5个图片（每个20KB），RTT=50ms，链路带宽为2Mbps。计算使用非持久连接和持久连接（流水线）下载该网页的总时间。
\tcblower
\textbf{传输时间计算}：
\begin{itemize}[itemsep=0pt]
    \item HTML传输时间 = $\frac{2\times8}{2000}$ = 8ms
    \item 每个图片传输时间 = $\frac{20\times8}{2000}$ = 80ms
\end{itemize}

\textbf{非持久连接}：每个对象需要2RTT + 传输时间
\begin{itemize}[itemsep=0pt]
    \item HTML：2×50 + 8 = 108ms
    \item 5个图片：5×(2×50 + 80) = 5×180 = 900ms
    \item 总时间 = 108 + 900 = \textbf{1008ms}
\end{itemize}

\textbf{持久连接（流水线）}：
\begin{itemize}[itemsep=0pt]
    \item 建立连接：1RTT = 50ms
    \item 请求HTML：1RTT + 8ms = 58ms
    \item 连续请求5个图片：1RTT + 5×80ms = 450ms
    \item 总时间 ≈ 50 + 58 + 450 = \textbf{558ms}
\end{itemize}
\end{examplebox}

\begin{examplebox}[简答题：DNS缓存作用]
\textbf{题目}：DNS缓存的作用是什么？缓存记录会永久保存吗？
\tcblower
\textbf{DNS缓存的作用}：
\begin{enumerate}[itemsep=0pt]
    \item 减少DNS查询时间，加快域名解析速度
    \item 减轻DNS服务器的负载
    \item 减少网络流量
\end{enumerate}

\textbf{缓存不会永久保存}：
\begin{itemize}[itemsep=0pt]
    \item 每条DNS记录都有TTL（生存时间）
    \item TTL到期后，缓存记录会被删除
    \item 下次查询时需要重新向DNS服务器请求
    \item 这样可以保证域名解析结果的及时更新
\end{itemize}
\end{examplebox}

\begin{examplebox}[简答题：Cookie与Session区别]
\textbf{题目}：比较Cookie和Session的区别。
\tcblower
\begin{center}\small
\begin{tabular}{|l|c|c|}
\hline
\textbf{特性} & \textbf{Cookie} & \textbf{Session} \\
\hline
存储位置 & 客户端 & 服务器端 \\
安全性 & 较低 & 较高 \\
存储容量 & 有限（约4KB） & 较大 \\
生命周期 & 可设置过期时间 & 会话结束或超时 \\
服务器压力 & 小 & 大 \\
跨域支持 & 不支持 & 不支持 \\
\hline
\end{tabular}
\end{center}

\textbf{联系}：Session通常需要Cookie来存储Session ID，以便服务器识别用户。
\end{examplebox}
